\documentclass[a4paper, 11pt, oneside]{book}

\usepackage{fontspec}
\usepackage[italian]{babel}
\usepackage{amsmath}
\usepackage{tikz}
\usepackage{XCharter}
\usepackage[xcharter]{newtxmath}
\usepackage[hidelinks]{hyperref}
\pagestyle{plain}
\usepackage[most]{tcolorbox}
\usepackage{float}
\usepackage{graphicx}
\usepackage{tabularx}
\usepackage{booktabs}
\usepackage{needspace}

\tcbuselibrary{skins, breakable, listings}
\setcounter{tocdepth}{1}
\setlength{\skip\footins}{25pt}
\usetikzlibrary{positioning, fit, arrows.meta, backgrounds}
\defaultfontfeatures{} % XCharter imposta Extension=.otf: va resettato
\setmonofont{JetBrainsMono Nerd Font}[
Scale=MatchLowercase
]

\definecolor{defborder}{HTML}{1F3FA8}
\definecolor{defback}{HTML}{F4F6FB}

\definecolor{notaborder}{HTML}{D97706}
\definecolor{notaback}{HTML}{FFFBEB}

\definecolor{codeborder}{HTML}{707070}
\definecolor{codeback}{HTML}{F0F0F0}

\newtcolorbox{definizione}[1][Definizione]{ breakable, enhanced, sharp corners, boxrule=0.6pt, colframe=defborder, colback=defback, title=#1, fonttitle=\bfseries, coltitle=defborder, attach title to upper={\par\smallskip}, before skip=12pt, after skip=12pt, boxsep=6pt }

\newtcolorbox{nota}[2]{ enhanced, sharp corners, boxrule=0.6pt, colframe=notaborder, colback=notaback, title={#1 -- #2}, fonttitle=\bfseries, coltitle=notaborder, attach title to upper={\par\smallskip}, before skip=12pt, after skip=12pt, boxsep=6pt }

\definecolor{codekeyword}{HTML}{1F3FA8}
\definecolor{codecomment}{HTML}{6A737D}
\definecolor{codestring}{HTML}{A31515}

\lstdefinelanguage{pseudo}{ morekeywords={function, end, returns, return, if, then, else, while, do, for, each, in, loop, repeat, until, persistent, static, local, and, or, not, true, false, null, failure}, sensitive=true, alsoletter={-}, morecomment=[l]{//}, morestring=[b]" }

\newtcblisting{codice}[1][Codice]{ breakable, enhanced, sharp corners, boxrule=0.6pt, colframe=codeborder, colback=codeback, title=#1, fonttitle=\bfseries, coltitle=codeborder, attach title to upper={\par\smallskip}, before skip=12pt, after skip=12pt, boxsep=4pt, left=4pt, right=4pt, grow to left by=1cm, grow to right by=1cm, listing only, listing options={ language=pseudo, basicstyle=\ttfamily\footnotesize, keywordstyle=\color{codekeyword}\bfseries, commentstyle=\color{codecomment}\itshape, stringstyle=\color{codestring}, breaklines=true, breakatwhitespace=true, breakautoindent=true, breakindent=2em, postbreak=\mbox{\textcolor{codecomment}{$\hookrightarrow$}\space}, tabsize=4, columns=fullflexible, keepspaces=true } }

\title{Intelligenza Artificiale}
\author{Michele Castenedoli}
\date{1 ottobre 2026}

\begin{document}
	\maketitle

	\tableofcontents

	\chapter{Introduzione all'Intelligenza Artificiale}
	\section{Definizione e Storia}
	Nel 1950 Alan Turing pubblica un articolo intitolato ``Computing Machinery and
	Intelligence'' in cui viene proposto un esperimento per determinare se una
	macchina pu\`o essere considerata dotata di intelligenza. Tale esperimento,
	noto come \emph{Test di Turing}, coinvolge un interrogatore umano che comunica
	con due entit\`a nascoste: una \textbf{macchina} e un \textbf{essere umano}. 
    L'interrogatore sottopone domande ad entrambe le entit\`a e, basandosi sulle risposte 
    fornite, deve determinare quale delle due \`e la macchina. Se l'interrogatore non \`e in
	grado di distinguere le due parti, la macchina \`e considerata intelligente.

	In seguito l'attenzione si \`e spostata sulla ricerca di metodi utili per
	risolvere problemi che richiedono l'intelligenza umana mediante l'utilizzo di
	algoritmi e modelli matematici fino ad arrivare all'introduzione di reti
	neurali e deep learning.

	\begin{definizione}
		[Intelligenza Artificiale] L'\textbf{Intelligenza Artificiale} \`e la disciplina
		che si occupa di studiare come \textbf{simulare} l'intelligenza umana in
		scenari complessi.
	\end{definizione}

	\section{Tipi di Intelligenza Artificiale}
	\subsection{Autonomous Agents}
	Gli \textbf{Autonomous Agents} sono sistemi che percepiscono l'ambiente e
	agiscono autonomamente per raggiungere determinati obiettivi.

	\subsection{Data Analysis}
	Consiste nell'utilizzo di algoritmi per analizzare grandi quantit\`a di dati
	ed estrarre informazioni utili e possibili correlazioni complesse.

	\subsection{Machine Learning}
	\`E lo sviluppo di algoritmi che permettono a modelli di apprendere da \emph{dati
	di esempio} e migliorare le loro prestazioni nel tempo senza essere
	esplicitamente programmati: a differenza della programmazione manuale e tradizionale,
	dove lo sviluppatore scrive una serie di istruzioni codificate in un
	linguaggio di programmazione, si fornisce soltanto un algoritmo di base e una grande
	quantit\`a di dati da cui apprendere. Il sistema analizza tali dati per
	trovare possibili pattern, correlazioni e regole statistiche necessarie per svolgere
	il compito.

	L'apprendimento automatico \`e suddiviso in tre categorie principali:
	\begin{itemize}
		\item \textbf{Unsupervised Learning}: il modello viene addestrato nel corso del
			tempo su un insieme di dati non etichettati (senza alcun ``tag'', categoria,
			o classificazione assegnati a priori da un essere umano). Il target \`e scoprire
			diverse strutture nascoste o pattern nei dati senza l'utilizzo di risposte
			corrette predefinite.

		\item \textbf{Supervised Learning}: il modello viene addestrato su un insieme
			di dati etichettati, dove ogni esempio di input \`e associato ad una risposta
			corretta. L'obiettivo \`e che il modello impari a mappare gli input alle risposte
			corrette.

		\item \textbf{Reinforcement Learning}: il modello impara attraverso continue
			interazioni con l'ambiente, ricevendo ricompense o penalit\`a a seconda
			delle azioni intraprese. L'obiettivo \`e massimizzare la ricompensa totale
			nel tempo. Pertanto il sistema valuta le conseguenze delle scelte a lungo termine
			anzich\'e cercare un vantaggio immediato.

			\emph{Esempio}: il sistema che simula un giocatore di scacchi potrebbe sacrificare
			la regina se prevede che in due mosse sar\`a in grado di dare scacco matto,
			raggiungendo cos\`i la massima ricompensa disponibile.
	\end{itemize}

	\subsection{Time Series Analysis}
	L'analisi delle serie temporali \`e un'area dell'apprendimento automatico che si
	concentra sull'analisi di dati collezionati nel tempo. Le serie temporali sono
	sequenze di dati misurati ad intervalli regolari (temperatura giornaliera,
	prezzi delle azioni o dati di vendita mensili). In questo caso, l'obiettivo \`e
	identificare pattern, tendenze, stagionalit\`a\footnote{Stagionalit\`a: la presenza
	di fluttuazioni, variazioni o pattern regolari prevedibili che si ripetono a
	intervalli regolari di tempo.} per fare previsioni future.

	Gli approcci comuni per l'analisi delle serie temporali sono:
	\begin{itemize}
		\item \textbf{Riconoscimento di anomalie e cause}: processo di identificazione
			di dati o eventi che si discostano significativamente dal comportamento normale
			o atteso. Queste anomalie possono fornire indicatori di problemi, errori o
			situazioni non previste che richiedono attenzione.

		\item \textbf{Generative Transformers}: sono una classe di modelli che permettono
			di predire il prossimo elemento di una sequenza di dati partendo dagli elementi
			precedenti (token di un testo, pixel di un'immagine). Si introduce il concetto
			di \textbf{attenzione}: meccanismo usato per pesare l'importanza delle diverse
			parti della sequenza in input durante la generazione di output.

			\emph{Esempio}: Se il modello di apprendimento sta elaborando la sequenza di
			token relativa alla frase ``ho mangiato una pesca'', il meccanismo alla
			base di attenzione assegner\`a un peso (punteggio) molto elevato alle
			parole ``mangiato'' e ``una'' per aiutare il sistema a capire che ci si
			sta riferendo al frutto e non all'hobby di pescare.
	\end{itemize}

	\subsection{Intelligent Agents}
	Un agente intelligente \`e un sistema che percepisce l'ambiente circostante
	attraverso sensori e agisce sull'ambiente per raggiungere un obiettivo
	specifico e svolgere un determinato compito. Gli elementi cardine di un agente
	intelligente includono:
	\begin{itemize}
		\item \textbf{Performance Measure}: misurazione del successo (efficacia) dell'agente
			nel raggiungere i suoi obiettivi.

		\item \textbf{Rationality}: l'agente deve agire in modo da massimizzare la sua
			performance measure attesa, dati la sequenza percettiva e la conoscenza a
			priori di cui dispone.
	\end{itemize}

	\begin{nota}
		{Differenza}{Intelligent Agents vs Autonomous Agents} I due concetti si
		sovrappongono in gran parte, ma pongono l'accento su aspetti diversi. 
        L'\textbf{intelligenza} riguarda la \emph{qualit\`a delle decisioni}: un 
        agente \`e intelligente (razionale) se sceglie le azioni che massimizzano 
        la performance measure attesa. L'\textbf{autonomia}
		riguarda invece \emph{quanto l'agente dipende dal progettista}: un agente \`e
		autonomo nella misura in cui il suo comportamento si basa sulla propria esperienza
		(percezioni e apprendimento) anzich\'e solo sulla conoscenza inserita a priori,
		e opera senza un intervento umano continuo. Le due propriet\`a sono quindi 
        indipendenti: un termostato con regole fisse \`e un agente ma non \`e autonomo, 
        mentre un'auto a guida autonoma che impara e si adatta a situazioni nuove \`e sia
		intelligente sia autonoma.
	\end{nota}

	\subsection{Markov Decision Process (MDP)}
	Un MDP \`e un modello matematico utilizzato per rappresentare problemi di decisione
	sequenziali in ambienti stocastici. Gli elementi principali sono:

	\begin{itemize}
		\item \textbf{States}: insieme $S$ degli stati; uno stato $s \in S$
			rappresenta l'ambiente in un dato momento

		\item \textbf{Actions}: insieme $A$ delle azioni che l'agente pu\`o intraprendere

		\item \textbf{Transition model}: effetto che le azioni hanno sull'ambiente.
			Poich\'e gli effetti potrebbero essere incerti, il modello assegna 
            una probabilit\`a a ogni possibile stato successivo
			\begin{displaymath}
				\text{T}(s, a, s') = P(s' \mid s, a)
			\end{displaymath}
			Dove $T$ \`e la funzione di transizione e esprime la probabilit\`a di
			capitare nello stato $s'$ sapendo che si \`e nello stato s e che l'azione $a$
			viene eseguita.

			\begin{nota}
				{Approfondimento}{Transition Model} Da notare che $T$ \`e solo un altro modo
				per descrivere la probabilit\`a condizionata di capitare in $s'$ sapendo
				che si parte dallo stato $s$ e si esegue l'azione $a$.
			\end{nota}

		\item \textbf{Reward}: valore \textbf{immediato} dell'esecuzione di un'azione
			\begin{displaymath}
				R: S \times A \times S \rightarrow \mathbb{R}
			\end{displaymath}
			Si definisce la ricompensa che l'agente guadagna raggiungendo quel passo:
			\begin{itemize}
				\item Il primo $S$ indica lo stato di partenza

				\item $A$ \`e l'azione che \`e stata eseguita

				\item Il secondo $S$ indica lo stato in cui si \`e giunti
			\end{itemize}

		\item \textbf{Policy}: strategia che l'agente utilizza per decidere quale azione
			intraprendere in ogni stato con l'obiettivo di massimizzare la ricompensa totale
			attesa nel tempo
			\begin{displaymath}
				\pi : S \rightarrow A
			\end{displaymath}
	\end{itemize}

	Alla base del modello c'\`e la \textbf{propriet\`a di Markov}, da cui prende
	il nome: la probabilit\`a di raggiungere lo stato successivo dipende \emph{soltanto}
	dallo stato corrente e dall'azione scelta, e non dalla storia degli stati e delle
	azioni precedenti:
	\begin{displaymath}
		P(s_{t+1}\mid s_{t}, a_{t}, s_{t-1}, a_{t-1}, \dots, s_{0}) = P(s_{t+1}\mid s
		_{t}, a_{t})
	\end{displaymath}
	In altre parole, lo stato corrente riassume tutta l'informazione rilevante per
	decidere come agire.

	\chapter{Agenti Intelligenti}
	\section{Definizione}
	\begin{definizione}
		[Agente Intelligente] Un \textbf{agente} \`e qualsiasi cosa possa essere vista
		come un sistema che percepisce il suo ambiente attraverso dei \textbf{sensori}
		e agisce su di esso tramite \textbf{attuatori}.
	\end{definizione}

	\emph{Esempio}: L'agente umano possiede sensori come occhi e orecchie, mentre
	pu\`o usare come attuatori mani, gambe, tratto vocale. Un agente robotico potrebbe
	invece avere telecamere e telemetri a infrarossi per sensori e diversi motori come
	attuatori.

	\vspace{2em}
% \begin{noindent}
	\begin{figure}[H]
		\centering
		\resizebox{0.67\textwidth}{!}{
		\begin{tikzpicture}[
			>=Stealth,
			font=\sffamily,
			thick,
			box/.style={draw, rounded corners=15pt, minimum height=7.5cm, fill=gray!10},
			blackbox/.style={draw, fill=black, text=white, minimum size=2.2cm, 
            font=\Huge\bfseries}
		]
			% 1. Nodi centrali
			\node[blackbox] (brain) {?};
			\node[above=1.2cm of brain, font=\large] (sensori) {sensori};
			\node[below=1.2cm of brain, font=\large] (attuatori) {attuatori};

			% Frecce interne
			\draw[->] (sensori.south) -- (brain.north);
			\draw[->] (brain.south) -- (attuatori.north);

			% 2. Box dell'Agente
			\begin{scope}[on background layer]
				\node[box, fit=(brain) (sensori) (attuatori), minimum width=6cm,
					inner sep=1cm] (agente) {};
			\end{scope}

			% Etichetta "Agente" centrata in alto con padding verticale (yshift)
			\node[below, font=\Large, yshift=-0.4cm] at (agente.north) {Agente};

			% 3. Box dell'Ambiente
			\node[box, right=3.5cm of agente, minimum width=2cm] (ambiente) {};
			\node[rotate=-90, font=\Large] at (ambiente.center) {Ambiente};

			% 4. Frecce esterne con etichette spostate a destra (pos=0.65)
			\draw[<-] (sensori.east) -- node[below, pos=0.65] {percezioni}
				(sensori.east -| ambiente.west);
			\draw[->] (attuatori.east) -- node[above, pos=0.65] {azioni}
				(attuatori.east -| ambiente.west);
		\end{tikzpicture}
		}
	\end{figure}
% \end{noindent}

	\vspace{1em}
	Si usa il termine \textbf{percezione} per indicare i dati che i sensori raccolgono.
	La \textbf{Sequenza Percettiva} di un agente \`e la storia di tutto ci\`o che
	l'agente ha raccolto in tutta la sua esistenza. In generale, la scelta dell'azione
	di un agente in un qualsiasi istante pu\`o dipendere dalla conoscenza
	integrata in esso e dall'intera sequenza percettiva osservata fino a quel
	momento, ma non da qualcosa che l'agente non abbia percepito.

	Specificando l'azione che un agente pu\`o scegliere in qualunque istante, si
	definisce l'agente in modo completo. In termini matematici, il comportamento
	di un agente \`e descritto dalla funzione agente, che descrive la
	corrispondenza tra una qualsiasi sequenza percettiva e una specifica azione.

	La funzione agente assume la seguente forma:
	\begin{displaymath}
		f: \mathcal{P}^{*}\rightarrow \mathcal{A}
	\end{displaymath}

	In particolare:
	\begin{itemize}
		\item $\mathcal{P}$ indica l'insieme delle singole percezioni

		\item $^{*}$ (\emph{chiusura di Kleene}) \`e l'operatore che, applicato a un
			insieme, restituisce l'insieme di \emph{tutte} le sequenze finite (compresa
			quella vuota) formate dai suoi elementi

		\item $\mathcal{P}^{*}$ \`e quindi l'insieme di tutte le possibili sequenze
			percettive: ogni suo elemento rappresenta una possibile cronologia delle percezioni
			registrate fino a un certo momento

		\item $\mathcal{A}$ indica l'insieme delle possibili azioni
	\end{itemize}

	Internamente, la funzione agente \`e implementata da un \textbf{programma
	agente}.
	\begin{nota}
		{Differenza}{Programma Agente vs Funzione Agente} Mentre la funzione agente
		\`e una descrizione matematica astratta, il programma agente \`e una sua 
        implementazione fisica concreta, in esecuzione all'interno del sistema fisico.
	\end{nota}
	\emph{Esempio}: Un mondo dell'aspirapolvere con due sole posizioni. Ogni posizione
	pu\`o essere pulita o sporca, l'agente pu\`o andare a sinistra o a destra e pu\`o
	pulire il riquadro in cui si trova.
	\begin{center}
		\begin{tabular}{ |c|c| }
			\hline
			Sequenza Percettiva                      & Azione   \\
			\hline
			{}[\emph{A, Pulito}]                     & Destra   \\
			{}[\emph{A, Sporco}]                     & Aspira   \\
			{}[\emph{B, Pulito}]                     & Sinistra \\
			{}[\emph{B, Sporco}]                     & Aspira   \\
			{}[\emph{A, Pulito}], [\emph{A, Pulito}] & Destra   \\
			{}[\emph{A, Pulito}], [\emph{A, Sporco}] & Aspira   \\
			$\vdots$                                 & $\vdots$ \\
			\hline
		\end{tabular}
	\end{center}
	\vspace{1em}
	Se un agente ha $\vert{}\mathcal{P}\vert{}$ possibili percezioni, dopo $T$
	passi temporali, la funzione agente ha $\sum_{t=1}^{T}\vert{}\mathcal{P}\vert{}
	^{t}$ voci. Se lo storico di percezioni \`e irrilevante, cio\`e se l'azione
	dipende soltanto dalla percezione corrente, l'agente viene detto \textbf{Simple
	Reflex Agent}.

	L'obiettivo \`e progettare piccoli programmi agente per rappresentare enormi funzioni
	agente.

% \begin{noindent}
\begin{codice}
[Un possibile programma agente]
function AGENTE-REATTIVO-ASPIRAPOLVERE([posizione, stato]) returns azione
    if stato = Sporco then return Aspira
    else if posizione = A then return Destra
    else if posizione = B then return Sinistra
\end{codice}
% \end{noindent}

	Questo programma \`e adeguato se la performance measure prende in considerazione
	solo il numero di stanze pulite. Se subentrassero altre metriche che portano
	ad una definizione della performance measure diversa (stato della batteria,
	percentuale di ricarica, minimizzare il numero di spostamenti), tale programma
	non \`e in grado di esprimere la cosa migliore da fare. Da notare che l'agente
	che implementa tale programma \`e un \textbf{Simple Reflex Agent}: l'azione
	dipende soltanto dalla percezione corrente.

	\section{Razionalit\`a}
	\subsection{Comportarsi Correttamente: il concetto di razionalit\`a}

	\begin{definizione}
		Un \textbf{agente razionale} \`e un agente che fa la cosa giusta.
	\end{definizione}

	Esistono diverse accezioni dell'espressione ``cosa giusta'', ma l'IA rimane fedele
	alla nozione di \textbf{consequenzialismo}: si valuta il comportamento di un agente
	considerando le conseguenze. Quando un agente viene inserito in un ambiente, genera
	una sequenza di azioni in base alle percezioni che riceve. Questa sequenza di
	azioni porta l'ambiente ad attraversare una sequenza di stati: se tale sequenza
	\`e desiderabile, significa che l'agente si \`e comportato bene. Questa nozione
	di desiderabilit\`a \`e catturata da una misura di prestazione che valuta la
	sequenza di stati dell'ambiente.

	In un dato momento, ci\`o che \`e razionale dipende da quattro fattori:
	\begin{itemize}
		\item la misura di prestazione che definisce il criterio del successo

		\item la conoscenza pregressa dell'ambiente da parte dell'agente

		\item le azioni che l'agente pu\`o eseguire

		\item la sequenza percettiva dell'agente fino allo stato corrente
	\end{itemize}

	Tutto ci\`o porta a una definizione molto formale di \textbf{agente razionale}:
	\begin{definizione}
		Per ogni possibile sequenza di percezioni, un agente razionale dovrebbe
		scegliere un'azione che massimizzi il valore atteso della sua misura di
		prestazione, date le informazioni fornite dalla sequenza percettiva e da ogni
		ulteriore conoscenza dell'agente.
	\end{definizione}

	La razionalit\`a allora non \`e altro che un problema di ottimizzazione sulla
	performance measure: un agente razionale sceglie l'azione che massimizza il valore
	atteso della misura di prestazione, data la sequenza percettiva accumulata
	fino a quel momento.

	Notare come essere razionale non vuol dire essere onnisciente.\footnote{Possedere
	conoscenza illimitata su qualsiasi cosa.} Infatti le percezioni potrebbero non
	fornire tutte le informazioni rilevanti in un dato istante. Pertanto, i
	risultati potrebbero essere diversi da quelli previsti.

	\begin{nota}
		{Differenza}{Razionalit\`a vs Perfezione} La razionalit\`a massimizza la
		prestazione \textbf{attesa}, mentre la perfezione massimizzerebbe la
		prestazione \textbf{effettiva}. Un agente razionale non pu\`o essere
		giudicato per esiti che non poteva prevedere: se attraverso la strada dopo aver
		controllato che sia libera e vengo colpito da un oggetto caduto dall'alto,
		la mia scelta \`e stata comunque razionale.
	\end{nota}

	\subsection{Esercizio}
	L'esercizio propone di analizzare un ambiente composto da tre stanze (A, B, C)
	e due robot ($r_{1}, r_{2}$). Il robot $r_{1}$ pu\`o pattugliare le stanze A e
	B partendo da A, mentre $r_{2}$ pu\`o pattugliare B e C partendo da C.
	Assumendo un tempo di spostamento nullo, l'obiettivo \`e definire un
	comportamento razionale capace di minimizzare l'inattivit\`a media (\emph{average
	idleness}), definita come la frazione del tempo in cui una stanza non \`e visitata
	da alcun robot.

	Per affrontare il problema, possiamo analizzare tre diverse strategie:
	\begin{itemize}
		\item \textbf{Prima strategia (Vettore di inattivit\`a $[0, 1, 0]$):} I
			robot non cambiano mai la loro posizione iniziale. Il robot $r_{1}$ rimane
			in A e $r_{2}$ in C. Di conseguenza, le stanze A e C hanno inattivit\`a nulla,
			ma la stanza B non viene mai coperta, generando un'inattivit\`a pari a 1.

		\item \textbf{Seconda strategia (Vettore di inattivit\`a $[1/2, 1/2, 0]$):}
			Si decide di muovere soltanto il primo robot ($r_{1}$), che si 
            sposter\`a continuamente da A a B e viceversa, mentre $r_{2}$ resta fermo
            in C. In questo scenario, le stanze A e B rimangono scoperte per met\`a del tempo.

		\item \textbf{Terza strategia (Vettore di inattivit\`a $[1/4, 1/2, 1/4]$):}
			Questa strategia consente di far giocare un ruolo attivo anche al secondo
			robot ($r_{2}$), coordinando il movimento di entrambi per coprire l'intera
			area: ogni robot trascorre $3/4$ del tempo nella propria stanza esterna e
			$1/4$ in B, e i due non si trovano mai in B contemporaneamente. La stanza
			B resta quindi coperta per met\`a del tempo.
	\end{itemize}

	Si noti che in tutte e tre le strategie la somma delle inattivit\`a vale 1, quindi
	l'\textbf{inattivit\`a media \`e sempre $1/3$}: rispetto all'obiettivo formale
	le tre strategie sono equivalenti. Tuttavia la \textbf{terza strategia \`e la migliore}:
	nella pratica, per coprire il pi\`u possibile tutte le diverse stanze in modo omogeneo,
	\`e fondamentale introdurre il concetto di \emph{equit\`a}. Per garantire questa
	equit\`a, limitarsi a valutare la media non \`e sufficiente: risulta
	necessario introdurre il calcolo della \textbf{varianza} o della \textbf{deviazione
	standard}, in modo da penalizzare gli squilibri estremi tra le stanze.

	\section{PEAS}
	Per progettare un agente intelligente bisogna definire l'ambiente in cui opera:
	\begin{itemize}
		\item \textbf{Performance Measures}: metodo di valutazione del successo dell'agente

		\item \textbf{Environment}: il contesto in cui l'agente opera

		\item \textbf{Actuators}: i mezzi attraverso cui l'agente agisce sull'ambiente

		\item \textbf{Sensors}: i mezzi attraverso cui l'agente percepisce l'ambiente
	\end{itemize}

	\needspace{8\baselineskip}
	\emph{Esempio}: Prendiamo in considerazione un taxi automatico: il PEAS
	potrebbe essere:
	\begin{table}[H]
		\centering
		\renewcommand{\arraystretch}{1.3}
		\begin{tabularx}{\textwidth}{@{}l X@{}}
			\toprule
			\textbf{Performance measures} & Soddisfazione del cliente, sicurezza, efficienza
				del carburante (SOLO E85), rispetto per le leggi stradali \\
			\textbf{Environment} & Traffico stradale, condizioni meteorologiche,
				segnali stradali, pedoni e altri veicoli \\
			\textbf{Actuators} & Volante, acceleratore, freni, indicatori di direzione \\
			\textbf{Sensors} & Telecamere, GPS, sensori di velocit\`a, radar \\
			\bottomrule
		\end{tabularx}
	\end{table}

	\section{Tipi di Ambiente}
	Gli ambienti possono essere classificati secondo diverse caratteristiche:
	\begin{itemize}
		\item \textbf{Osservabile} (vs \emph{parzialmente osservabile}): se l'agente
			pu\`o percepire completamente lo stato dell'ambiente in ogni momento

		\item \textbf{Deterministico} (vs \emph{stocastico}): se lo stato successivo
			dell'ambiente \`e determinato completamente dallo stato corrente e dall'azione
			eseguita dall'agente

		\item \textbf{Episodico} (vs \emph{sequenziale}): se l'esperienza dell'agente
			\`e divisa in episodi indipendenti, cio\`e l'azione in un episodio non
			influisce sugli episodi successivi

		\item \textbf{Statico} (vs \emph{dinamico}): se l'ambiente non cambia mentre
			l'agente sta prendendo una decisione

		\item \textbf{Discreto} (vs \emph{continuo}): se stati, percezioni, azioni e
			tempo assumono un numero finito di valori distinti

		\item \textbf{Singolo Agente} (vs \emph{multi-agente}): se l'agente opera da
			solo nell'ambiente senza la presenza di altri agenti

		\item \textbf{Noto} (vs \emph{ignoto}): se l'agente (o il progettista)
			conosce le regole dell'ambiente, cio\`e gli esiti delle azioni
	\end{itemize}
	\begin{nota}
		{Approfondimento}{Scenario Singolo Agente in contesti Multi-Agenti} Un
		ambiente multi-agente pu\`o essere considerato singolo agente in quelle
		circostanze in cui si prevede la progettazione di un singolo agente mentre
		gli altri agenti non vengono modellati (che quindi vengono considerati elementi
		dell'ambiente).
	\end{nota}

	\emph{Esempio}: Tipi di ambienti in diversi contesti
	\begin{table}[h]
		\centering
		\begin{tabular}{|l|c|c|c|c|}
			\hline
			                                   & Crossword & Robo-selector & Poker  & Taxi   \\
			\hline
			\textcolor{magenta}{Observable}    & Yes       & Partly        & Partly & Partly \\
			\hline
			\textcolor{magenta}{Deterministic} & Yes       & No            & No     & No     \\
			\hline
			\textcolor{magenta}{Episodic}      & No        & Yes           & No     & No     \\
			\hline
			\textcolor{magenta}{Static}        & Yes       & No            & Yes    & No     \\
			\hline
			\textcolor{magenta}{Discrete}      & Yes       & No            & Yes    & No     \\
			\hline
			\textcolor{magenta}{Single-agent}  & Yes       & Yes           & No     & No     \\
			\hline
		\end{tabular}
	\end{table}

	Durante il corso ci soffermeremo sugli ambienti \emph{osservabili} e \emph{deterministici}
	per descrivere problemi sequenziali: un'azione svolta ha un impatto sul futuro.

	\section{Tipi di Problemi}
	\begin{itemize}
		\item \textbf{Deterministico, completamente osservabile} $\implies$ \textbf{problema
            a stato singolo (Single-State Problem)}. \\L'agente sa esattamente in quale 
            stato si trover\`a; la soluzione \`e una sequenza di azioni.

		\item \textbf{Non osservabile} $\implies$ \textbf{problema conformante}. \\L'agente
			potrebbe non avere idea di dove si trovi: l'agente non dispone di sensori.
            la soluzione (se esiste) \`e una sequenza di azioni.

		\item \textbf{Non deterministico e/o parzialmente osservabile} $\implies$
			\textbf{problema di contingenza}. \\Le percezioni forniscono costantemente
			nuove informazioni sullo stato corrente; la soluzione \`e un piano
			contingente (basato su condizioni del tipo if-then)
            o una \textit{policy} (mapping stato-azione): molto spesso si alternano 
            le fasi di ricerca e di esecuzione.

		\item \textbf{Spazio degli stati sconosciuto} $\implies$ \textbf{problema di
			esplorazione} (\textit{online}). \\L'agente deve scoprire gli stati e le
			transizioni interagendo direttamente con l'ambiente per capire come funziona 
            il problema.
	\end{itemize}

	\section{Programma Agente}
	Poich\'e l'agente \`e composto da un'architettura interna e un \textbf{programma
	agente}, \`e utile definire le caratteristiche generali di un programma agente.

	In generale, lo scheletro di un programma \`e composto da:
	\begin{itemize}
		\item Input: Percezione \textbf{corrente}. In tutti quei contesti in cui l'agente
			necessita di una sequenza di percezioni, ogni percezione dovr\`a essere salvata
			per riassemblare l'intera sequenza

		\item Output: \textbf{Prossima} azione
	\end{itemize}

	In questo schema \`e possibile seguire alcune architetture che consentono la modellazione
	di diversi agenti. Architetture diverse hanno diversi scopi. Di seguito l'elenco
	dei quattro tipi base di quasi tutti i sistemi intelligenti:
	\begin{itemize}
		\item Simple Reflex Agents

		\item Reflex Agents with State

		\item Agenti basati su obiettivi (Goal-Based Agents)

		\item Agenti basati sull'utilit\`a (Utility-Based Agents)
	\end{itemize}

	\subsection{Simple Reflex Agents}
	L'architettura pi\`u semplice \`e l'\textbf{Agente Reattivo Semplice}. Questi agenti
	scelgono le azioni sulla base della percezione corrente, ignorando tutta la
	storia percettiva precedente. Pertanto, il principio seguito \`e quello del
	tipo \textbf{regola condizione-azione}.
% \begin{noindent}
	\begin{figure}[H]
		\centering
		\begin{tikzpicture}[
			font=\sffamily\small,
			>={Stealth[length=2.5mm]},
			box/.style={draw, thick, align=center, minimum width=3.2cm, minimum height=1cm},
			regole/.style={draw, thick, rounded corners=8pt, 
            minimum height=0.6cm, inner xsep=8pt}
		]
			% Elementi interni all'agente
			\node (sensori)   at (0,0)                    {sensori};
			\node[box] (mondo)     [below=0.6cm of sensori]    {aspetto corrente\\ del mondo};
			\node[box] (azione)    [below=1.8cm of mondo]      {azione da\\ eseguire adesso};
			\node[regole] (reg)    [left=1.0cm of azione]      {regole condizione-azione};
			\node  (attuatori)     [below=0.6cm of azione]     {attuatori};

			% Riquadro agente
			\draw[thick, rounded corners=15pt] (-7,0.6) rectangle (2,-6.3);
			\node[anchor=north west, font=\sffamily\large] at (-6.8,0.45) {agente};

			% Riquadro ambiente
			\draw[thick, rounded corners=15pt] (2.4,0.6) rectangle (4.4,-6.3);
			\node[rotate=-90, font=\sffamily\large] at (3.4,-2.85) {ambiente};

			% Frecce
			\draw[->, thick] (3.9,0) -- (sensori.east);
			\draw[->, thick] (sensori) -- (mondo);
			\draw[->, thick] (mondo) -- (azione);
			\draw[->, thick] (reg) -- (azione);
			\draw[->, thick] (azione) -- (attuatori);
			\draw[->, thick] (attuatori.east) -- (3.4,0 |- attuatori.east);

			% Cornice esterna
			\draw (-7.4,1) rectangle (4.8,-6.7);
		\end{tikzpicture}
	\end{figure}
% \end{noindent}

% \begin{noindent}
\begin{codice}
[Simple Reflex Agents]
function AGENTE-REATTIVO-SEMPLICE(percezione) returns azione
    persistent: regole      // insieme di regole condizione-azione

    stato ← INTERPRETA-INPUT(percezione)
    regola ← REGOLA-CORRISPONDENTE(stato, regole)
    azione ← regola.AZIONE
    return azione
\end{codice}
% \end{noindent}

	Gli agenti reattivi semplici hanno il pregio della semplicit\`a, ma un'intelligenza
	molto limitata: funzionano correttamente \textbf{solo se l'ambiente \`e completamente
	osservabile}, cio\`e se la decisione giusta pu\`o essere presa sulla base
	della sola percezione corrente. Anche una minima parte di non osservabilit\`a pu\`o
	causare seri problemi.

	\emph{Esempio}: se l'aspirapolvere non avesse il sensore di posizione e percepisse
	soltanto [\emph{Pulito}] o [\emph{Sporco}], davanti a una stanza pulita
	dovrebbe scegliere sempre la stessa azione, ad esempio \emph{Sinistra}. Se si trova
	in A, quell'azione fallisce sempre: l'agente entra in un \textbf{ciclo
	infinito}.

	\subsection{Reflex Agents with State}
	Il modo pi\`u efficace di gestire l'osservabilit\`a parziale, per un agente, \`e
	tener traccia della parte del mondo che non pu\`o vedere nell'istante corrente.
	Questo significa che l'agente deve mantenere una sorta di \textbf{stato
	interno} che dipende dalla storia delle percezioni e che quindi riflette
	almeno una parte degli aspetti non osservabili dello stato corrente. In particolare,
	aggiornare lo stato interno richiede due componenti fondamentali: gli effetti delle
	azioni dell'agente e le modalit\`a di evoluzione del mondo indipendentemente dall'agente.
	Tale conoscenza viene implementata tramite un \textbf{modello di transizione}.
	Inoltre, sono necessari \textbf{modelli sensoriali} che descrivono come lo
	stato del mondo si rifletta nelle percezioni dell'agente.

	Un agente che utilizza tali modelli prende il nome di \textbf{Agente Basato su
	Modello}.
% \begin{noindent}
	\begin{figure}[H]
		\centering
		\begin{tikzpicture}[
			font=\sffamily\small,
			>={Stealth[length=2.5mm]},
			box/.style={draw, thick, align=center, minimum width=3.2cm, minimum height=1cm},
			tondo/.style={draw, thick, rounded corners=8pt, minimum height=0.6cm, 
            inner xsep=8pt}
		]
			% Elementi interni all'agente
			\node (sensori)   at (0,0)                    {sensori};
			\node[box] (mondo)     [below=0.6cm of sensori]    {aspetto corrente\\ del mondo};
			\node[box](azione)    [below=2cm of mondo]        {azione da\\ eseguire adesso};
			\node (attuatori) [below=0.6cm of azione]     {attuatori};

			% Elementi con bordi smussati
			\node[tondo, anchor=east] (stato)  at (-3.2,-0.5) {stato};
			\node[tondo, anchor=east] (evolve) at (-2.6,-1.3) {come evolve il mondo};
			\node[tondo, anchor=east] (cosa)   at (-2.6,-2.4){cosa fanno le mie azioni};
			\node[tondo, anchor=east] (regole) at ([xshift=-1cm]azione.west) 
            {regole condizione-azione};

			% Riquadro agente
			\draw[thick, rounded corners=15pt] (-7,0.6) rectangle (2,-6.3);
			\node[anchor=south west, font=\sffamily\large] at (-6.6,-6.1) {agente};

			% Riquadro ambiente
			\draw[thick, rounded corners=15pt] (2.4,0.6) rectangle (4.4,-6.3);
			\node[rotate=-90, font=\sffamily\large] at (3.4,-2.85) {ambiente};

			% Frecce
			\draw[->, thick] (3.9,0) -- (sensori.east);
			\draw[->, thick] (sensori) -- (mondo);
			\draw[->, thick] (stato.east) -- ([yshift=2mm]mondo.west);
			\draw[->, thick] (evolve.east) -- ([yshift=-1mm]mondo.west);
			\draw[->, thick] (cosa.east) -- ([xshift=3mm]mondo.south west);
			\draw[->, thick, dashed] ([xshift=6mm]mondo.north west)
				to[out=100, in=20] (stato.north east);
			\draw[->, thick] (mondo) -- (azione);
			\draw[->, thick] (regole.east) -- (azione.west);
			\draw[->, thick] (azione) -- (attuatori);
			\draw[->, thick] (attuatori.east) -- (3.4,0 |- attuatori.east);

			% Cornice esterna
			\draw (-7.4,1) rectangle (4.8,-6.7);
		\end{tikzpicture}
	\end{figure}
% \end{noindent}

% \begin{noindent}
\begin{codice}
[Reflex Agents with State]
function AGENTE-REATTIVO-BASATO-SU-MODELLO(percezione) returns azione
    persistent:
        stato               // concezione corrente dello stato del mondo
        modello_transizione // stato successivo in funzione di stato e azione
        modello_sensoriale  // come lo stato si riflette nelle percezioni
        regole              // insieme di regole condizione-azione
        azione              // l'azione più recente, inizialmente nessuna

    stato ← AGGIORNA-STATO(stato, azione, percezione,
                           modello_transizione, modello_sensoriale)
    regola ← REGOLA-CORRISPONDENTE(stato, regole)
    azione ← regola.AZIONE
    return azione
\end{codice}
% \end{noindent}

	\subsection{Goal-Based Agents}
	La conoscenza dello stato corrente dell'ambiente non \`e sempre sufficiente per
	decidere quale azione intraprendere. Spesso \`e necessario introdurre informazioni
	relative all'\textbf{obiettivo} (\textit{goal}), ovvero la descrizione delle
	situazioni desiderabili da raggiungere. Un agente basato su obiettivi (\textit{Goal-Based
	Agent}) combina queste informazioni con il proprio modello interno del mondo, superando
	il limite delle semplici regole ``condizione-azione''. Non basta pi\`u, infatti,
	associare direttamente una percezione a un'azione prestabilita. A differenza degli
	agenti reattivi (\textit{Reflex Agents}), in cui l'obiettivo \`e solo
	implicito nelle regole pre-programmate, l'agente basato su obiettivi valuta in
	modo esplicito le conseguenze delle proprie scelte: prima di agire, analizza attivamente
	se lo svolgimento di una determinata azione lo porter\`a o meno al soddisfacimento
	del proprio scopo.

	\emph{Esempio}: L'agente reattivo frena quando vede le luci di frenata, senza
	rendersi conto del motivo. Un agente basato su obiettivi frena quando vede le luci
	di frenata perch\'e quella \`e l'unica azione che, secondo la sua previsione,
	permette di raggiungere l'obiettivo di non tamponare altre macchine.

	\textbf{Ricerca} e \textbf{pianificazione} sono sottocampi dell'IA dedicati a
	identificare le sequenze di azioni da intraprendere per raggiungere gli
	obiettivi.

	Pu\`o sembrare che un Goal-Based Agent sia meno efficiente (magari l'agente che
	modella le funzionalit\`a di guida autonoma potrebbe frenare invece di accelerare
	su una strada dritta ad alta velocit\`a), ma garantisce maggiore \textbf{flessibilit\`a},
	perch\'e la conoscenza che guida le sue decisioni \`e rappresentata
	esplicitamente e pu\`o essere modificata.

	\subsection{Agenti Basati Sull'Utilit\`a}
	Raggiungere semplicemente un obiettivo non basta: per un risultato ottimale \`e
	fondamentale valutare anche la qualit\`a e l'efficienza delle azioni scelte per
	ottenerlo (come velocit\`a, sicurezza o convenienza).

	In questo ambito, infatti, si introduce il concetto di \textbf{Utilit\`a}: all'interno
	dell'agente viene implementata una \textbf{Funzione di utilit\`a} che internalizza
	la performance measure. Purch\'e la funzione di utilit\`a interna e la misura di
	prestazione esterna concordino, un agente che sceglie le azioni per massimizzare
	l'utilit\`a sar\`a comunque razionale in base alla misura di prestazione esterna.

	Si sottolinea che questo non \`e l'unico modo per essere razionali: l'esempio 
    sull'aspirapolvere mostra la modellazione di un comportamento razionale senza 
    introdurre obiettivi o funzioni di utilit\`a. Ma anche in questo caso i vantaggi sono
	principalmente in termini di flessibilit\`a, soprattutto in contesti in cui \`e
	necessario prendere decisioni in condizioni d'incertezza.

	In termini tecnici, un agente razionale basato sull'utilit\`a sceglie l'azione
	che massimizza l'\textbf{utilit\`a attesa} dei risultati, ovvero l'utilit\`a
	che l'agente si attende di ottenere, in media, date le probabilit\`a e le
	utilit\`a di ciascun risultato.

	\emph{Esempio}: Si immagini un taxi a guida autonoma che deve scegliere tra
	due strade:
	\begin{itemize}
		\item \textbf{STRADA A (Urbana)}: \`e sempre libera ($100\%$) e porta a
			destinazione nei tempi previsti, con utilit\`a moderata pari a 5. \\Utilit\`a
			attesa: $1.0 \times 5 = 5$

		\item \textbf{STRADA B (Autostrada)}: \`e libera nell'$80\%$ dei casi, offrendo
			la possibilit\`a di arrivare in largo anticipo, con utilit\`a elevata pari
			a 10. Nel restante $20\%$ dei casi, invece, si trova un ingorgo totale e
			si arriva molto in ritardo, con utilit\`a = $-5$. \\Utilit\`a attesa:
			$(0.8 \times 10) + (0.2 \times -5) = 8 - 1 = 7$
	\end{itemize}

	La strada che massimizza l'utilit\`a attesa \`e pertanto la B.

	\begin{nota}
		{Approfondimento}{Utilit\`a Attesa} Si noti che non \`e altro che il calcolo
		del valore medio (\emph{speranza matematica}): $E[X] = \sum x_{i} \cdot P(X =
		x_{i})$
	\end{nota}

	\begin{nota}
		{Differenza}{Funzione di Utilit\`a vs Performance Measure} Sebbene
		rappresentino concetti assai simili, la Performance Measure \`e il metro di
		giudizio \textbf{esterno}, stabilito dal progettista per valutare la qualit\`a
		complessiva del comportamento dell'agente: l'agente non ha necessariamente accesso
		diretto a essa.

		La funzione di utilit\`a \`e invece \textbf{interna} all'agente: \`e 
        l'internalizzazione della performance measure, costruita dal progettista in 
        modo che le due siano allineate. L'agente la usa per assegnare un 
        valore a ogni possibile stato del mondo e, in condizioni di incertezza, 
        sceglie l'azione che massimizza l'utilit\`a
		attesa. Se la funzione di utilit\`a non rispecchia fedelmente la performance
		measure, l'agente pu\`o massimizzarla e comportarsi comunque in modo non razionale
		rispetto al criterio esterno.
	\end{nota}

% \begin{noindent}
	\begin{figure}[H]
		\centering
		\resizebox{0.87\textwidth}{!}{
		\begin{tikzpicture}[
			font=\sffamily\small,
			>={Stealth[length=2.5mm]},
			box/.style={draw, thick, align=center, minimum width=3.9cm, minimum height=1cm},
			tondo/.style={draw, thick, rounded corners=8pt, minimum height=0.6cm, 
            inner xsep=8pt}
		]
			% Elementi interni all'agente
			\node (sensori)     at (0,0.2)                     {sensori};
			\node[box] (mondo)       [below=0.3cm of sensori]  {aspetto corrente\\ del mondo};
			\node[box] (conseguenze) [below=0.3cm of mondo] 
            {come cambier\`a il mondo \\ se eseguo l'azione A};
			
            \node[box] (utile)       [below=0.6cm of conseguenze]   
            {quanto sar\`o contento \\ in uno stato simile};
			
            \node[box] (azione)      [below=0.6cm of utile]         
            {azione da \\ eseguire adesso};
			
            \node (attuatori)        [below=0.4cm of azione]        {attuatori};

			% Fondo dei riquadri, calcolato in base agli attuatori
			\coordinate (fondo) at ([yshift=-3mm]attuatori.south);

			% Elementi con bordi smussati
			\node[tondo, anchor=east] (stato)  at (-3.2,-0.5) {stato};
			\node[tondo, anchor=east] (evolve) at (-2.6,-1.3) {come evolve il mondo};
			\node[tondo, anchor=east] (cosa)   at (-2.6,-2.4) {cosa fanno le mie azioni};
			\node[tondo, anchor=east] (regole) at ([xshift=-1cm]utile.west) {utilit\`a};

			% Frecce della colonna centrale
			\draw[->, thick] (3.9,0.2) -- (sensori.east);
			\draw[->, thick] (sensori) -- (mondo);
			\draw[->, thick] (mondo) -- (conseguenze);
			\draw[->, thick] (conseguenze) -- (utile);
			\draw[->, thick] (utile) -- (azione);
			\draw[->, thick] (azione) -- (attuatori);
			\draw[->, thick] (attuatori.east) -- (3.4,0 |- attuatori.east);

			% Frecce dagli elementi con bordi smussati
			\draw[->, thick, dashed] ([xshift=6mm]mondo.north west)
				to[out=100, in=20] (stato.north east);
			\draw[->, thick] (stato.east) -- ([yshift=2mm]mondo.west);
			\draw[->, thick] (evolve.east) -- ([yshift=-1mm]mondo.west);
			\draw[->, thick] (evolve.east) -- ([yshift=3mm]conseguenze.west);
			\draw[->, thick] (cosa.east) -- ([yshift=-2mm]mondo.west);
			\draw[->, thick] (cosa.east) -- (conseguenze.west);
			\draw[->, thick] (regole.east) -- (utile.west);

			% Riquadro agente
			\draw[thick, rounded corners=15pt] (-7,0.6) rectangle (2.2,0 |- fondo);
			\node[anchor=north west, font=\sffamily\large] at (-6.8,0.45) {agente};

			% Riquadro ambiente
			\draw[thick, rounded corners=15pt] (2.4,0.6) rectangle (4.4,0 |- fondo);
			\path (3.4,0.6) -- (3.4,0 |- fondo)
				node[midway, rotate=-90, font=\sffamily\large] {ambiente};

			% Cornice esterna
			\draw (-7.4,1) rectangle ([yshift=-4mm]4.8,0 |- fondo);
		\end{tikzpicture}
		}
	\end{figure}
% \end{noindent}

	\chapter{Risolvere Problemi con la Ricerca}

	\section{Problem Solving Agents}
	Quando l'azione da svolgere non \`e subito evidente, un agente potrebbe avere
	la necessit\`a di guardare avanti, cio\`e considerare una sequenza di azioni 
    (\emph{cammino})
	che porta a uno stato obiettivo. Questo tipo di agente viene chiamato \textbf{Agente
	Risolutore di Problemi} (\emph{Problem Solving Agent}). Gli agenti risolutori di
	problemi utilizzano rappresentazioni \textbf{atomiche} degli stati, i quali vengono
	considerati come entit\`a prive di una struttura interna.

	I Problem Solving Agents sono in particolare una sottocategoria dei Goal-Based
	Agents. Per decidere quale azione intraprendere, questi agenti adottano un approccio
	sistematico sequenziale basato sulla formulazione di problemi, determinazione
	di obiettivi e sulla ricerca di soluzioni. Tale processo si divide in quattro macro-fasi:
	\begin{enumerate}
		\item \textbf{Formulazione dell'obiettivo (\emph{Goal Formulation})}: dato lo
			stato corrente dell'ambiente, l'agente definisce una situazione
			desiderabile (l'obiettivo) da raggiungere.

		\item \textbf{Formulazione del problema (\emph{Problem Formulation})}: l'agente
			decide quali stati e quali azioni prendere in considerazione per
			raggiungere tale obiettivo.

		\item \textbf{Ricerca (\emph{Search})}: l'agente esplora lo spazio degli stati
			simulando diverse sequenze di azioni nel proprio modello interno, fino a
			trovare una soluzione valida (una sequenza di mosse che porta all'obiettivo).

		\item \textbf{Esecuzione (\emph{Execution})}: l'agente esegue, una alla
			volta, le azioni previste dalla soluzione trovata, ignorando le percezioni.
	\end{enumerate}
	Notare come la risoluzione del problema in questo caso \`e di tipo \textbf{offline}:
	la soluzione viene eseguita a ``occhi chiusi'' (come detto prima, si ignorano le
	percezioni).

	\begin{nota}
		{Differenza}{Online vs Offline Problem Solving} A differenza della tecnica
		offline, la quale pu\`o essere usata in ambienti completamente deterministici
		e osservabili, la \textbf{Online Problem Solving} prevede che l'agente agisca
		in condizioni di parziale ignoranza, dove la conoscenza sull'ambiente non \`e
		completa. Pertanto \`e indispensabile che i sensori rimangano attivi per valutare
		l'esito delle proprie azioni.
	\end{nota}

    Di seguito un'implementazione di un programma agente per un Problem Solving Agent
    che procede 
    alla risoluzione di un problema mediante ricerca nello spazio degli stati.

% \begin{noindent}
\begin{codice}
[Simple Problem-Solving Agent]
function AGENTE-RISOLUTORE-SEMPLICE(percezione) returns azione
    persistent:
        seq         // sequenza di azioni, inizialmente vuota
        stato       // descrizione dello stato corrente del mondo
        obiettivo   // un obiettivo, inizialmente nullo
        problema    // la formulazione del problema

    stato ← AGGIORNA-STATO(stato, percezione)
    if seq è vuota then // la prima volta che si esegue
        obiettivo ← FORMULA-OBIETTIVO(stato) // che si vuole raggiungere
        problema ← FORMULA-PROBLEMA(stato, obiettivo) 
        seq ← RICERCA(problema)
        if seq = failure then return null
    azione ← PRIMO(seq)
    seq ← RESTO(seq)
    return azione
\end{codice}
% \end{noindent}

    Da notare come l'agente sceglie cosa fare \textbf{solamente} all'inizio e poi 
    continua, ignorando completamente le percezioni.

    \section{Come Formulare un Problema}
    \subsection{Problema: Vacanze in Romania (con un M5 e39)}
    Ci troviamo in Romania e dobbiamo andare a Bucarest per poi prendere l'aereo 
    e tornare a casa.
    Attualmente ci troviamo ad Arad. Vogliamo trovare la serie di citt\`a da attraversare per 
    giungere a Bucarest.
    Dobbiamo individuare un corretto livello di astrazione, ma ragionevolmente si procede nel
    seguente modo:
    \begin{itemize}
        \item \textbf{Obiettivo}: andare a Bucarest
        \item \textbf{Formulazione del problema}
            \begin{enumerate}
                \item \textbf{states}: le diverse citt\`a
                \item \textbf{actions}: guidare per attraversare le citt\`a 
            \end{enumerate}
        \item \textbf{Soluzione da Individuare}: sequenza di citt\`a da attraversare per 
            arrivare a Bucarest
    \end{itemize}

    Fin da subito \`e possibile mettere in discussione il livello di deterministicit\`a 
    del problema. A causa di agenti esterni (incidenti, lavori) si potrebbe 
    essere dirottati su citt\`a non scelte. Tali supposizioni vengono inizialmente ignorate. 
    In questo caso vogliamo trovare \textbf{una} soluzione: non \`e detto che quella 
    trovata sia a costo minimo.
    %TODO: importare immagine grafo topologico ROMANIA
    \begin{center}
        \includegraphics[width=0.7\textwidth]{images/romania_mappa}
    \end{center}

    Cerchiamo di proporre una formulazione di tale quesito mediante Single-State problem. 
    Allora si individuano i seguenti elementi:
    \begin{enumerate}
        \item \textbf{Stato Iniziale}: la citt\`a di partenza.

        \item \textbf{Successor Function} che assume la seguente forma: 
            $f : S \times A \rightarrow S$.
            Si mappa uno stato e un'azione in un altro stato. 
            \emph{Esempio}: $S(A) = \{<Arad \rightarrow Zerind, Zerind>, \hdots \}$

        \item \textbf{Goal Test} che pu\`o essere di due nature diverse:
            \begin{enumerate}
                \item esplicito: la citt\`a di destinazione (Bucarest). Lo stato in cui 
                    ci si deve trovare affinch\'e il problema possa considerarsi risolto.

                \item implicito: una feature che descrive pi\`u di uno stato.
            \end{enumerate}

        \item \textbf{Path Cost}: Non sempre si desidera trovare una soluzione qualunque. 
            Potrebbe essere utile trovare quella a costo minimo. Nel nostro esempio 
            corrisponde ad una metrica che indica il costo di un percorso 
            (il costo espresso sugli archi). Il path cost \`e un processo additivo in 
            cui vengono sommati gli \textbf{step cost}, 
            funzioni del tipo $STEP\_COST = S \times A \times S \rightarrow \mathbb{R}^+$.
        \item \textbf{Soluzione}: sequenza di azioni che porta dallo stato iniziale a 
            uno dei goal-state.
    \end{enumerate}

    \section{Tree-Search Algorithm}
    L'idea di base \`e fare un'esplorazione simulata e offline 
    (essendo un single state problem)
    le percezioni vengono ignorate. Concettualmente l'agente ragiona, individua una sequenza di
    azioni e procede a eseguirla. Pertanto si generano dei successori degli stati usando 
    regole che si conoscono del mondo.

% \begin{noindent}
\begin{codice}
[Tree-Search Algorithm]
function TREE-SEARCH(problema, strategia) returns una soluzione, o failure
    inizializza l'albero di ricerca usando lo stato iniziale del problema

    loop do
        if nessun candidato per l'espansione then return failure
        scelgo uno dei nodi per l'espansione seguendo la strategia
        if nodo contiene un goal-state then return soluzione
        else aggiungi il nodo successore al search tree (e espando)
    end
\end{codice}
% \end{noindent}

	\subsection{Stati vs Nodi}
	Nell'albero di ricerca \`e importante distinguere due concetti che a prima vista
	possono sembrare equivalenti:
	\begin{itemize}
		\item uno \textbf{stato} \`e la rappresentazione di una configurazione fisica del
			mondo: descrive \emph{dove} si trova l'agente (ad esempio, ``sono ad Arad'');

		\item un \textbf{nodo} \`e una \textbf{struttura dati} che costituisce parte
			dell'albero di ricerca: oltre allo stato a cui corrisponde, memorizza
			\emph{come} ci si \`e arrivati.
	\end{itemize}

	I campi di un nodo sono:
	\begin{table}[H]
		\centering
		\renewcommand{\arraystretch}{1.3}
		\begin{tabularx}{\textwidth}{@{}l X@{}}
			\toprule
			\textbf{State} & lo stato del mondo a cui corrisponde il nodo \\
			\textbf{Parent} & il nodo da cui \`e stato generato \\
			\textbf{Action} & l'azione applicata al padre per generare il nodo \\
			\textbf{Children} & i nodi generati espandendo il nodo \\
			\textbf{Depth} & il numero di passi dalla radice \\
			\textbf{Path cost} $g(x)$ & il costo del cammino dalla radice fino al nodo \\
			\bottomrule
		\end{tabularx}
	\end{table}

	\begin{nota}
		{Differenza}{Stato vs Nodo}
		Gli stati \textbf{non} hanno padre, azioni, figli, profondit\`a o costo di
		cammino: queste informazioni dipendono dal \emph{percorso} seguito e non dal
		mondo. Per questo lo \textbf{stesso stato} pu\`o comparire in \textbf{pi\`u
		nodi diversi} dell'albero di ricerca.
	\end{nota}

	\emph{Esempio}: nel problema della Romania, lo stato ``sono a Sibiu'' pu\`o essere
	raggiunto:
	\begin{itemize}
		\item direttamente da Arad: profondit\`a 1, $g = 140$;

		\item passando per Arad $\rightarrow$ Zerind $\rightarrow$ Oradea $\rightarrow$
			Sibiu: \\ profondit\`a 3, ${g = 75 + 71 + 151 = 297}$.
	\end{itemize}
	Si ottengono quindi due nodi distinti, con padre, profondit\`a e costo diversi, che
	contengono per\`o lo stesso stato. In sintesi: lo stato indica \emph{dove} si trova
	l'agente, il nodo indica anche \emph{come} ci \`e arrivato.

\end{document}
